\section{Introduction}

\noindent\textcolor{titleblue}{\Large\bfseries L}'accès à l'information juridique est un droit
fondamental dans tout État de droit. Au Maroc, le corpus juridique est vaste et
complexe, comprenant des centaines de lois, codes et décrets rédigés
principalement en arabe classique moderne (MSA). Pour les citoyens, les
praticiens du droit et les chercheurs, naviguer dans cet océan de textes
juridiques afin de trouver une réponse précise à une question spécifique
constitue un défi majeur. Les moteurs de recherche généralistes (Google, Bing)
sont mal adaptés à cette tâche : ils ne comprennent ni la spécificité du
vocabulaire juridique arabe ni la structure hiérarchique des codes de loi
(articles, chapitres, sections).

\vspace{0.2cm}

\noindent Les récents progrès en Traitement du Langage Naturel (NLP),
particulièrement dans le domaine des modèles de langage de grande taille (LLM)
et du Retrieval-Augmented Generation (RAG), offrent une opportunité sans
précédent pour construire des systèmes de Question-Réponse (Q\&R) spécialisés
\cite{lewis2020retrieval}. L'approche RAG combine un module de recherche
d'information (retrieval) avec un module de génération de texte (LLM), permettant
de produire des réponses contextualisées, sourcées et à jour sans nécessiter de
fine-tuning coûteux. Cependant, l'application du RAG au domaine juridique arabe
présente des défis uniques : la rareté des modèles d'embedding multilingues
optimisés pour le vocabulaire juridique, la complexité morphologique de l'arabe
\cite{habash2010arabic}, la présence de Darija (arabe dialectal marocain) dans
les requêtes des utilisateurs, et la nécessité d'une fidélité absolue aux textes
de loi. Les documents sources sont eux-mêmes difficiles d'accès : la majorité
des textes historiques existe uniquement au format PDF scanné, nécessitant une
chaîne complète d'OCR et de correction par modèle de vision (VLM) avant toute
exploitation.

\vspace{0.2cm}

\noindent Le système que nous proposons s'appuie sur un corpus de \textbf{2\,357
articles juridiques} couvrant sept domaines du droit marocain : pénal (822
articles), trafic (228), civil (705), commercial (359), organisation judiciaire
(32), fiscal (211), et contraventions. Ces articles sont extraits de leurs PDF
sources, nettoyés, normalisés en arabe standard, structurés avec métadonnées
hiérarchiques, et indexés dans PostgreSQL + pgvector à l'aide du modèle
d'embedding \texttt{intfloat/multilingual-e5-large} (1024 dimensions).

\vspace{0.2cm}

\noindent Cet article présente la conception, l'implémentation et l'évaluation
rigoureuse d'un système complet de Q\&R juridique marocain basé sur RAG. Notre
contribution est quadruple :

\begin{enumerate}
    \item Une \textbf{chaîne de traitement NLP arabe complète} allant de
    l'extraction de texte à partir de PDF (numériques et scannés) jusqu'à la
    normalisation morphologique, le nettoyage, et l'analyse structurelle des
    codes de loi marocains.
    \item Une \textbf{architecture de recherche hybride} combinant similarité
    vectorielle dense (modèle \texttt{BAAI/bge-m3}, 1024 dimensions) avec
    recherche lexicale BM25, fusionnée par Reciprocal Rank Fusion (RRF), et
    indexée dans PostgreSQL + pgvector avec un index HNSW.
    \item Un \textbf{benchmark systématique} — le premier à notre connaissance
    — de cinq modèles d'embedding multilingues (\texttt{intfloat/multilingual-e5-large},
    \texttt{BAAI/bge-m3}, \texttt{LaBSE}, \texttt{mpnet-base}, \texttt{e5-base})
    sur un corpus de 26 articles juridiques marocains avec 65 variantes de
    requêtes couvrant l'arabe formel, le Darija, le français, l'anglais, et
    les requêtes courtes/longues.
    \item Un \textbf{benchmark comparatif} de trois stratégies de recherche
    (dense, BM25, RRF) sur 32 questions réparties sur 7 domaines juridiques,
    évaluant à la fois la qualité de la récupération (Précision@5, Rappel@5,
    MRR@5, NDCG@5) et la qualité des réponses générées (pertinence, fidélité,
    exactitude).
\end{enumerate}

\vspace{0.2cm}

\noindent La suite de cet article est organisée comme suit : la section 2
présente les travaux connexes en NLP arabe, systèmes RAG juridiques, et
modèles d'embedding multilingues. La section 3 décrit l'architecture globale
du système. La section 4 détaille la méthodologie incluant le pipeline de
prétraitement, la stratégie de recherche hybride, et la génération de réponses.
La section 5 présente les résultats expérimentaux des benchmarks. La section 6
discute les implications, limitations, et directions futures. Enfin, la section
7 conclut l'article.